\section{End-Term Examination --- Set B}
\label{sec:endterm_set_b}

\LPUPaperHeader{End-Term Examination}{B}{3 Hours}{60 Marks}{CO1 to CO6 (Units 1 to 6 Comprehensive)}

\begin{center}
\sffamily\textbf{\color{AlertRed}Structure of Question Paper:}\\[1mm]
\textbf{Part A (Objective):} 30 MCQs strictly from \textbf{Units 4, 5, and 6}. Total weight: \textbf{20 Marks}. All questions compulsory.\\
\textbf{Part B (Subjective):} 6 Questions from \textbf{all units} (10 Marks each). \textbf{Attempt any 4 questions} ($4 \times 10 = \mathbf{40\text{ Marks}}$).\\
\textbf{Total Maximum Marks: 60 Marks.}
\end{center}

\vspace{3mm}

% ------------------------------------------------------------------------------
% PART A: 30 MCQS (UNITS 4, 5, 6)
% ------------------------------------------------------------------------------
\subsection*{PART A: Objective Section (30 MCQs --- Units 4, 5, 6) [20 Marks]}

\mcqitem{1}{The function $f(x,y) = \frac{x^2 - y^2}{x^2 + y^2}$ at $(0,0)$ has:}{Limit equal to $0$}{Limit equal to $1$}{No limit (limit does not exist)}{Limit equal to $-1$}{1}

\mcqitem{2}{If $z = e^{x/y}$, then $x \frac{\partial z}{\partial x} + y \frac{\partial z}{\partial y}$ is:}{$0$}{$z$}{$2z$}{$-z$}{1}

\mcqitem{3}{The value of the Jacobian $\frac{\partial(u,v)}{\partial(x,y)}$ when $u = x^2 - y^2$ and $v = 2xy$ is:}{$4(x^2 + y^2)$}{$2(x^2 + y^2)$}{$x^2 + y^2$}{$4xy$}{1}

\mcqitem{4}{The maximum value of the directional derivative of $\phi$ at any point is along the:}{Tangent vector}{Gradient vector $\nabla \phi$}{Normal to the gradient}{Coordinate axis}{1}

\mcqitem{5}{If $\vec{F} = (x+3y)\hat{i} + (y-2z)\hat{j} + (x+az)\hat{k}$ is solenoidal, then the constant $a$ is:}{$-2$}{$-1$}{$2$}{$0$}{1}

\mcqitem{6}{The scalar potential $\phi$ for the conservative force $\vec{F} = 2xy\hat{i} + x^2\hat{j}$ is:}{$x^2 y + C$}{$2x y^2 + C$}{$x^2 + y^2 + C$}{$xy + C$}{1}

\mcqitem{7}{The value of $\iint_R dx \, dy$ where $R$ is the ellipse $\frac{x^2}{a^2} + \frac{y^2}{b^2} \le 1$ is:}{$\pi ab$}{$4\pi ab$}{$\frac{\pi ab}{4}$}{$2\pi ab$}{1}

\mcqitem{8}{The value of $\int_0^a \int_0^{\sqrt{a^2-x^2}} dy \, dx$ is:}{$\frac{\pi a^2}{4}$}{$\pi a^2$}{$\frac{\pi a^2}{2}$}{$a^2$}{1}

\mcqitem{9}{If $\vec{r} = x\hat{i} + y\hat{j} + z\hat{k}$ and $r = \|\vec{r}\|$, then $\nabla(1/r)$ is:}{$-\frac{\vec{r}}{r^3}$}{$\frac{\vec{r}}{r^3}$}{$-\frac{\vec{r}}{r^2}$}{$\frac{1}{r^2}$}{1}

\mcqitem{10}{A vector field $\vec{F}$ is conservative if and only if:}{$\nabla \times \vec{F} = \vec{0}$}{$\nabla \cdot \vec{F} = 0$}{$\nabla^2 \vec{F} = 0$}{$\vec{F} = 0$}{1}

\mcqitem{11}{The value of $\int_0^1 \int_0^1 \int_0^1 (x + y + z) \, dx \, dy \, dz$ is:}{$3/2$}{$1$}{$3$}{$1/2$}{1}

\mcqitem{12}{The stationary points of $f(x,y) = x^2 + y^2$ are:}{$(0,0)$ only}{$(1,1)$ only}{$(0,1)$ and $(1,0)$}{Infinitely many}{1}

\mcqitem{13}{For $f(x,y) = xy$, the origin $(0,0)$ is a:}{Local Minimum}{Local Maximum}{Saddle Point}{Inconclusive}{1}

\mcqitem{14}{The value of $\nabla \times (\nabla \phi)$ for any scalar field $\phi$ is:}{$\vec{0}$}{$1$}{$\nabla^2 \phi$}{$\vec{r}$}{1}

\mcqitem{15}{The value of the line integral $\oint_C (x dy - y dx)$ along the circle $x^2 + y^2 = 1$ in counter-clockwise direction is:}{$2\pi$}{$\pi$}{$0$}{$4\pi$}{1}

\mcqitem{16}{Stokes' Theorem equates $\oint_C \vec{F} \cdot d\vec{r}$ to:}{$\iint_S (\nabla \times \vec{F}) \cdot \hat{n} \, dS$}{$\iiint_V (\nabla \cdot \vec{F}) \, dV$}{$\iint_S \vec{F} \cdot \hat{n} \, dS$}{$\oint_C (\nabla \cdot \vec{F}) \, ds$}{1}

\mcqitem{17}{The value of $\nabla \cdot (\nabla \times \vec{A})$ for any vector field $\vec{A}$ is:}{$0$}{$1$}{$\nabla^2 \vec{A}$}{$\vec{A}$}{1}

\mcqitem{18}{The cylindrical coordinates $(\rho, \theta, z)$ of the Cartesian point $(1, 1, 1)$ are:}{$(\sqrt{2}, \pi/4, 1)$}{$(2, \pi/4, 1)$}{$(\sqrt{2}, \pi/2, 1)$}{$(1, \pi/4, 1)$}{1}

\mcqitem{19}{The double integral $\int_0^\infty \int_0^\infty e^{-(x^2+y^2)} \, dx \, dy$ evaluates to:}{$\frac{\pi}{4}$}{$\pi$}{$\frac{\pi}{2}$}{$\sqrt{\pi}$}{1}

\mcqitem{20}{The angle between the surfaces $x^2 + y^2 + z^2 = 9$ and $z = x^2 + y^2 - 3$ at $(2, -1, 2)$ is found using their:}{Gradients}{Divergences}{Curls}{Laplacians}{1}

\mcqitem{21}{If $f(x,y) = x^3 y + y^3 x$, then $\frac{\partial^2 f}{\partial x \partial y}$ is:}{$3x^2 + 3y^2$}{$6xy$}{$x^3 + y^3$}{$3x^2 y^2$}{1}

\mcqitem{22}{The value of $\iint_S dS$ over the surface of a sphere of radius $a$ is:}{$4\pi a^2$}{$\frac{4}{3}\pi a^3$}{$2\pi a^2$}{$\pi a^2$}{1}

\mcqitem{23}{If $\vec{F} = x\hat{i} + y\hat{j}$, then $\nabla \cdot \vec{F}$ is:}{$2$}{$0$}{$1$}{$3$}{1}

\mcqitem{24}{The curve $C$ in Stokes' Theorem must be:}{A closed boundary curve of surface $S$}{An open straight line}{Infinite}{Discontinuous}{1}

\mcqitem{25}{The value of $\lim_{(x,y)\to(0,0)} \frac{x^3 + y^3}{x^2 + y^2}$ is:}{$0$}{$1$}{Does not exist}{$\infty$}{1}

\mcqitem{26}{The work done by a force $\vec{F}$ along a closed path in a conservative field is:}{$0$}{$1$}{Path dependent}{$\infty$}{1}

\mcqitem{27}{The Jacobian of the inverse transformation satisfies:}{$J \cdot J' = 1$}{$J + J' = 0$}{$J = J'$}{$J \cdot J' = 0$}{1}

\mcqitem{28}{If $u = f(x/y)$, then $x u_x + y u_y$ is:}{$0$}{$u$}{$x+y$}{$1$}{1}

\mcqitem{29}{The volume enclosed by the cylinder $x^2 + y^2 = 1$ between $z = 0$ and $z = 5$ is:}{$5\pi$}{$\pi$}{$10\pi$}{$25\pi$}{1}

\mcqitem{30}{The Laplacian operator $\nabla^2 \phi$ is defined as:}{$\nabla \cdot (\nabla \phi)$}{$\nabla \times (\nabla \phi)$}{$\nabla \phi \cdot \vec{r}$}{$\|\nabla \phi\|$}{1}

\vspace{5mm}
\hrule
\vspace{5mm}

% ------------------------------------------------------------------------------
% PART B: 6 SUBJECTIVE QUESTIONS (ATTEMPT ANY 4)
% ------------------------------------------------------------------------------
\subsection*{PART B: Descriptive Section (Attempt Any 4 of 6) [40 Marks]}

\questionitem{31 (Unit 1)}{
    Verify the \textbf{Cayley-Hamilton Theorem} for the matrix:
    \[
    A = \begin{bmatrix} 1 & 0 & 3 \\ 2 & 1 & -1 \\ 1 & -1 & 1 \end{bmatrix}
    \]
    and hence compute $A^{-1}$ and $A^4$.
}{10}{CO1}{L3 Apply}

\questionitem{32 (Unit 2)}{
    Solve the second-order non-homogeneous differential equation:
    \[
    (D^2 - 4D + 3)y = 20 \cos 2x + e^x
    \]
}{10}{CO2}{L3 Apply}

\questionitem{33 (Unit 3)}{
    \textbf{(a)} Find the radius of convergence and interval of convergence of the power series:
    \[
    \sum_{n=1}^\infty \frac{(-1)^n (x-2)^n}{n \cdot 3^n}
    \]
    \textbf{(b)} Test the convergence of the series $\sum_{n=1}^\infty \left(1 + \frac{1}{\sqrt{n}}\right)^{-n^{3/2}}$ using \textbf{Cauchy's Root Test}.
}{10}{CO3}{L4 Analyze}

\questionitem{34 (Unit 4)}{
    Find all the extreme values (maxima, minima, and saddle points) of the multivariable function:
    \[
    f(x,y) = x^4 + y^4 - 2(x - y)^2
    \]
}{10}{CO4}{L3 Apply}

\questionitem{35 (Unit 5)}{
    Evaluate the double integral by transforming to polar coordinates:
    \[
    \iint_R \frac{dx \, dy}{(1 + x^2 + y^2)^2}
    \]
    where $R$ is the region in the \textbf{first quadrant} bounded by the circle $x^2 + y^2 = 1$.
}{10}{CO5}{L3 Apply}

\questionitem{36 (Unit 6)}{
    Verify \textbf{Stokes' Theorem} for the vector field:
    \[
    \vec{F} = (2x - y)\hat{i} - y z^2\hat{j} - y^2 z\hat{k}
    \]
    where $S$ is the upper half surface of the sphere $x^2 + y^2 + z^2 = 1$ ($z \ge 0$) and its boundary curve $C$ is the circle $x^2 + y^2 = 1$ in the $xy$-plane.
}{10}{CO6}{L4 Analyze}

\vspace{5mm}
\hrule
\vspace{5mm}

% ------------------------------------------------------------------------------
% COMPLETE STEP-BY-STEP SOLUTIONS
% ------------------------------------------------------------------------------
\subsection*{Solutions \& Marking Scheme (End-Term Set B)}

\subsubsection*{Part A: Answer Key \& Explanations}

\begin{center}
\begin{tabular}{|c|c||c|c||c|c||c|c||c|c||c|c|}
\hline
\textbf{Q} & \textbf{Ans} & \textbf{Q} & \textbf{Ans} & \textbf{Q} & \textbf{Ans} & \textbf{Q} & \textbf{Ans} & \textbf{Q} & \textbf{Ans} & \textbf{Q} & \textbf{Ans} \\
\hline
1 & \textbf{C} & 6 & \textbf{A} & 11 & \textbf{A} & 16 & \textbf{A} & 21 & \textbf{A} & 26 & \textbf{A} \\
2 & \textbf{A} & 7 & \textbf{A} & 12 & \textbf{A} & 17 & \textbf{A} & 22 & \textbf{A} & 27 & \textbf{A} \\
3 & \textbf{A} & 8 & \textbf{A} & 13 & \textbf{C} & 18 & \textbf{A} & 23 & \textbf{A} & 28 & \textbf{A} \\
4 & \textbf{B} & 9 & \textbf{A} & 14 & \textbf{A} & 19 & \textbf{A} & 24 & \textbf{A} & 29 & \textbf{A} \\
5 & \textbf{A} & 10 & \textbf{A} & 15 & \textbf{A} & 20 & \textbf{A} & 25 & \textbf{A} & 30 & \textbf{A} \\
\hline
\end{tabular}
\end{center}

\begin{solutionbox}[Explanations for Critical MCQs]
\textbf{Q.1 (Ans: C)} Path $y = mx \implies \frac{1-m^2}{1+m^2}$ depends on slope $m$. Limit does not exist.\\[1.5mm]
\textbf{Q.3 (Ans: A)} $J = \begin{vmatrix} 2x & -2y \\ 2y & 2x \end{vmatrix} = 4x^2 - (-4y^2) = 4(x^2 + y^2)$.\\[1.5mm]
\textbf{Q.5 (Ans: A)} Solenoidal $\implies \nabla \cdot \vec{F} = 1 + 1 + a = 0 \implies a = -2$.\\[1.5mm]
\textbf{Q.15 (Ans: A)} By Green's Theorem, $\oint_C (x dy - y dx) = \iint_R (1 - (-1)) dA = 2 \iint_R dA = 2(\pi \cdot 1^2) = 2\pi$.\\[1.5mm]
\textbf{Q.19 (Ans: A)} $\int_0^{\pi/2} d\theta \int_0^\infty r e^{-r^2} dr = \frac{\pi}{2} [-\frac{1}{2}e^{-r^2}]_0^\infty = \frac{\pi}{4}$.
\end{solutionbox}

\subsubsection*{Part B: Detailed Subjective Solutions}

% Solution 31
\begin{solutionbox}[Solution to Question 31 (10 Marks)]
\textbf{Step 1: Characteristic Equation of $A$}
$A = \begin{bmatrix} 1 & 0 & 3 \\ 2 & 1 & -1 \\ 1 & -1 & 1 \end{bmatrix}$.
\begin{itemize}[leftmargin=6mm]
    \item $S_1 = 1 + 1 + 1 = 3$.
    \item $S_2 = M_{11} + M_{22} + M_{33} = (1 - 1) + (1 - 3) + (1 - 0) = 0 - 2 + 1 = -1$.
    \item $S_3 = \det(A) = 1(0) - 0 + 3(-2 - 1) = -9$.
\end{itemize}
Equation: $\lambda^3 - 3\lambda^2 - \lambda + 9 = 0$.

\textbf{Step 2: Verification of Cayley-Hamilton ($A^3 - 3A^2 - A + 9I = O$)}
Compute $A^2$:
\[
A^2 = \begin{bmatrix} 1 & 0 & 3 \\ 2 & 1 & -1 \\ 1 & -1 & 1 \end{bmatrix} \begin{bmatrix} 1 & 0 & 3 \\ 2 & 1 & -1 \\ 1 & -1 & 1 \end{bmatrix} = \begin{bmatrix} 4 & -3 & 6 \\ 3 & 2 & 4 \\ 0 & -2 & 5 \end{bmatrix}
\]
Compute $A^3 = A \cdot A^2$:
\[
A^3 = \begin{bmatrix} 1 & 0 & 3 \\ 2 & 1 & -1 \\ 1 & -1 & 1 \end{bmatrix} \begin{bmatrix} 4 & -3 & 6 \\ 3 & 2 & 4 \\ 0 & -2 & 5 \end{bmatrix} = \begin{bmatrix} 4 & -9 & 21 \\ 11 & -2 & 11 \\ 1 & -7 & 7 \end{bmatrix}
\]
Substitute into the polynomial:
\[
A^3 - 3A^2 - A + 9I = \begin{bmatrix} 4 & -9 & 21 \\ 11 & -2 & 11 \\ 1 & -7 & 7 \end{bmatrix} - 3\begin{bmatrix} 4 & -3 & 6 \\ 3 & 2 & 4 \\ 0 & -2 & 5 \end{bmatrix} - \begin{bmatrix} 1 & 0 & 3 \\ 2 & 1 & -1 \\ 1 & -1 & 1 \end{bmatrix} + 9\begin{bmatrix} 1 & 0 & 0 \\ 0 & 1 & 0 \\ 0 & 0 & 1 \end{bmatrix}
\]
\[
= \begin{bmatrix} 4 - 12 - 1 + 9 & -9 + 9 - 0 + 0 & 21 - 18 - 3 + 0 \\ 11 - 9 - 2 + 0 & -2 - 6 - 1 + 9 & 11 - 12 + 1 + 0 \\ 1 - 0 - 1 + 0 & -7 + 6 + 1 + 0 & 7 - 15 - 1 + 9 \end{bmatrix} = \begin{bmatrix} 0 & 0 & 0 \\ 0 & 0 & 0 \\ 0 & 0 & 0 \end{bmatrix} = O
\]
Hence, Cayley-Hamilton Theorem is **verified!**

\textbf{Step 3: Calculating $A^{-1}$}
Multiply by $A^{-1}$: $A^2 - 3A - I + 9A^{-1} = O \implies 9A^{-1} = -A^2 + 3A + I$:
\[
9A^{-1} = -\begin{bmatrix} 4 & -3 & 6 \\ 3 & 2 & 4 \\ 0 & -2 & 5 \end{bmatrix} + \begin{bmatrix} 3 & 0 & 9 \\ 6 & 3 & -3 \\ 3 & -3 & 3 \end{bmatrix} + \begin{bmatrix} 1 & 0 & 0 \\ 0 & 1 & 0 \\ 0 & 0 & 1 \end{bmatrix} = \begin{bmatrix} 0 & 3 & 3 \\ 3 & 2 & -7 \\ 3 & -1 & -1 \end{bmatrix}
\]
\[
\mathbf{A^{-1} = \frac{1}{9}\begin{bmatrix} 0 & 3 & 3 \\ 3 & 2 & -7 \\ 3 & -1 & -1 \end{bmatrix}}
\]

\begin{markingrubric}
\textbf{Mark Distribution:}
$\bullet$ Characteristic equation $\lambda^3 - 3\lambda^2 - \lambda + 9 = 0$: \textbf{3 Marks}\\
$\bullet$ Computation of $A^2, A^3$ and verification $= O$: \textbf{4 Marks}\\
$\bullet$ Exact calculation of $A^{-1}$: \textbf{3 Marks}
\end{markingrubric}
\end{solutionbox}

% Solution 32
\begin{solutionbox}[Solution to Question 32 (10 Marks)]
\textbf{Step 1: Complementary Function}
$m^2 - 4m + 3 = 0 \implies (m - 1)(m - 3) = 0 \implies m = 1, 3$.
\[
\mathbf{y_c = c_1 e^x + c_2 e^{3x}}
\]
\textbf{Step 2: Particular Integral for $e^x$ ($y_{p_1}$)}
Since $m = 1$ is a single root of the auxiliary equation:
\[
y_{p_1} = \frac{1}{(D-1)(D-3)} e^x = \frac{1}{(1-3)} \frac{1}{D-1} e^x = -\frac{1}{2} (x e^x) = \mathbf{-\frac{1}{2} x e^x}
\]
\textbf{Step 3: Particular Integral for $20\cos 2x$ ($y_{p_2}$)}
\[
y_{p_2} = \frac{1}{D^2 - 4D + 3} (20 \cos 2x)
\]
Replace $D^2$ with $-2^2 = -4$:
\[
= \frac{20}{-4 - 4D + 3} \cos 2x = \frac{20}{-4D - 1} \cos 2x = -\frac{20}{4D + 1} \cos 2x
\]
Multiply numerator and denominator by $(4D - 1)$:
\[
= -20 \frac{4D - 1}{16D^2 - 1} \cos 2x = -20 \frac{4D - 1}{16(-4) - 1} \cos 2x = -20 \frac{4D - 1}{-65} \cos 2x = \frac{4}{13}(4D - 1)\cos 2x
\]
Evaluate $D(\cos 2x) = -2\sin 2x$:
\[
= \frac{4}{13}[4(-2\sin 2x) - \cos 2x] = \mathbf{-\frac{4}{13}(8\sin 2x + \cos 2x)}
\]
\textbf{Step 4: Complete General Solution}
\[
\mathbf{y(x) = c_1 e^x + c_2 e^{3x} - \frac{1}{2} x e^x - \frac{32}{13}\sin 2x - \frac{4}{13}\cos 2x}
\]

\begin{markingrubric}
\textbf{Mark Distribution:}
$\bullet$ Complementary function $y_c = c_1 e^x + c_2 e^{3x}$: \textbf{2.5 Marks}\\
$\bullet$ Particular integral $y_{p_1}$ with case of failure: \textbf{3.5 Marks}\\
$\bullet$ Particular integral $y_{p_2}$ for trigonometric term: \textbf{3 Marks}\\
$\bullet$ Final general solution: \textbf{1 Mark}
\end{markingrubric}
\end{solutionbox}

% Solution 33
\begin{solutionbox}[Solution to Question 33 (10 Marks)]
\textbf{Part (a): Radius \& Interval of Convergence (5 Marks)}
Given power series: $\sum_{n=1}^\infty c_n (x-2)^n$ where $c_n = \frac{(-1)^n}{n 3^n}$.
\[
\frac{1}{R} = \lim_{n\to\infty} \left|\frac{c_{n+1}}{c_n}\right| = \lim_{n\to\infty} \frac{n \cdot 3^n}{(n+1) \cdot 3^{n+1}} = \frac{1}{3} \implies \mathbf{R = 3}
\]
The series converges for $|x - 2| < 3 \implies -3 < x - 2 < 3 \implies -1 < x < 5$.
\textbf{Testing Endpoints:}
\begin{itemize}[leftmargin=6mm]
    \item At $x = 5$: $\sum_{n=1}^\infty \frac{(-1)^n (3)^n}{n 3^n} = \sum_{n=1}^\infty \frac{(-1)^n}{n}$ (Alternating harmonic series). Converges by Leibniz test.
    \item At $x = -1$: $\sum_{n=1}^\infty \frac{(-1)^n (-3)^n}{n 3^n} = \sum_{n=1}^\infty \frac{(-1)^n (-1)^n 3^n}{n 3^n} = \sum_{n=1}^\infty \frac{1}{n}$ (Harmonic series). Diverges!
\end{itemize}
Hence, the interval of convergence is $\mathbf{(-1, 5]}$.

\textbf{Part (b): Cauchy Root Test (5 Marks)}
$u_n = \left(1 + \frac{1}{\sqrt{n}}\right)^{-n^{3/2}}$.
Apply Cauchy's Root Test:
\[
L = \lim_{n\to\infty} (u_n)^{1/n} = \lim_{n\to\infty} \left[ \left(1 + \frac{1}{\sqrt{n}}\right)^{-n^{3/2}} \right]^{1/n} = \lim_{n\to\infty} \left(1 + \frac{1}{\sqrt{n}}\right)^{-n^{1/2}} = \lim_{n\to\infty} \left(1 + \frac{1}{\sqrt{n}}\right)^{-\sqrt{n}}
\]
Let $k = \sqrt{n}$. As $n \to \infty, k \to \infty$:
\[
L = \lim_{k\to\infty} \left(1 + \frac{1}{k}\right)^{-k} = \frac{1}{e}
\]
Since $L = \frac{1}{e} \approx 0.368 < 1$, by Cauchy's Root Test, the series is \textbf{CONVERGENT}.

\begin{markingrubric}
\textbf{Mark Distribution:}
$\bullet$ Part (a) Radius $R = 3$: \textbf{2 Marks}\\
$\bullet$ Part (a) Endpoint analysis giving interval $(-1, 5]$: \textbf{3 Marks}\\
$\bullet$ Part (b) Root test power manipulation: \textbf{2.5 Marks}\\
$\bullet$ Part (b) Limit evaluation $1/e < 1$ and conclusion: \textbf{2.5 Marks}
\end{markingrubric}
\end{solutionbox}

% Solution 34
\begin{solutionbox}[Solution to Question 34 (10 Marks)]
\textbf{Step 1: First Partial Derivatives}
$f(x,y) = x^4 + y^4 - 2(x - y)^2 = x^4 + y^4 - 2x^2 + 4xy - 2y^2$.
\[
f_x = 4x^3 - 4x + 4y = 4(x^3 - x + y) = 0 \implies x^3 - x + y = 0 \label{eq:ext_b1}
\]
\[
f_y = 4y^3 + 4x - 4y = 4(y^3 + x - y) = 0 \implies y^3 + x - y = 0 \label{eq:ext_b2}
\]
\textbf{Step 2: Finding Critical Points}
Add \eqref{eq:ext_b1} and \eqref{eq:ext_b2}:
\[
x^3 + y^3 = 0 \implies y = -x
\]
Substitute $y = -x$ into \eqref{eq:ext_b1}:
\[
x^3 - x - x = 0 \implies x^3 - 2x = 0 \implies x(x^2 - 2) = 0
\]
This yields three real solutions:
\begin{itemize}[leftmargin=6mm]
    \item $x = 0 \implies y = 0 \implies \mathbf{P_1(0,0)}$.
    \item $x = \sqrt{2} \implies y = -\sqrt{2} \implies \mathbf{P_2(\sqrt{2}, -\sqrt{2})}$.
    \item $x = -\sqrt{2} \implies y = \sqrt{2} \implies \mathbf{P_3(-\sqrt{2}, \sqrt{2})}$.
\end{itemize}

\textbf{Step 3: Second Derivative Test}
\[
r = f_{xx} = 12x^2 - 4, \quad s = f_{xy} = 4, \quad t = f_{yy} = 12y^2 - 4
\]
Discriminant: $D = rt - s^2 = (12x^2 - 4)(12y^2 - 4) - 16$.
\begin{enumerate}[label=(\roman*)]
    \item At $P_1(0,0)$:
    $r = -4, t = -4, s = 4 \implies D = (-4)(-4) - 16 = 0$.
    Since $D = 0$, the test is inconclusive. Examine behavior along paths:
    Along $y = x$: $f(x,x) = 2x^4 \ge 0$ (positive near origin).
    Along $y = 0$: $f(x,0) = x^4 - 2x^2 = x^2(x^2 - 2) < 0$ for small $x$.
    Since $f$ takes both positive and negative values in every neighborhood of $(0,0)$, $\mathbf{(0,0)}$ is a \textbf{Saddle Point}.
    \item At $P_2(\sqrt{2}, -\sqrt{2})$:
    $r = 12(2) - 4 = 20 > 0$, $t = 12(2) - 4 = 20$.
    $D = (20)(20) - 16 = 384 > 0$.
    Since $D > 0$ and $r > 0$, this is a point of \textbf{Local Minimum}.
    Value: $f(\sqrt{2}, -\sqrt{2}) = 4 + 4 - 2(2\sqrt{2})^2 = 8 - 16 = \mathbf{-8}$.
    \item At $P_3(-\sqrt{2}, \sqrt{2})$:
    $r = 20 > 0, t = 20, D = 384 > 0$.
    This is also a point of \textbf{Local Minimum}.
    Value: $f(-\sqrt{2}, \sqrt{2}) = \mathbf{-8}$.
\end{enumerate}

\begin{markingrubric}
\textbf{Mark Distribution:}
$\bullet$ First partial derivatives and critical points $(0,0), (\pm\sqrt{2}, \mp\sqrt{2})$: \textbf{3.5 Marks}\\
$\bullet$ Second derivative test discriminant: \textbf{2.5 Marks}\\
$\bullet$ Investigating $(0,0)$ saddle point: \textbf{2 Marks}\\
$\bullet$ Classifying local minima at $(\pm\sqrt{2}, \mp\sqrt{2})$ with value $-8$: \textbf{2 Marks}
\end{markingrubric}
\end{solutionbox}

% Solution 35
\begin{solutionbox}[Solution to Question 35 (10 Marks)]
\textbf{Step 1: Polar Coordinate Transformation}
Let $x = r\cos\theta, y = r\sin\theta \implies x^2 + y^2 = r^2$, and $dx \, dy = r \, dr \, d\theta$.
The region $R$ is the first quadrant of the unit disk $x^2 + y^2 \le 1$:
\[
r \in [0, 1], \quad \theta \in \left[0, \frac{\pi}{2}\right]
\]
\textbf{Step 2: Transformed Integral}
\[
I = \int_0^{\pi/2} \int_0^1 \frac{r \, dr \, d\theta}{(1 + r^2)^2} = \left(\int_0^{\pi/2} d\theta\right) \left(\int_0^1 \frac{r \, dr}{(1 + r^2)^2}\right)
\]
\textbf{Step 3: Evaluation}
Angular integral: $\int_0^{\pi/2} d\theta = \frac{\pi}{2}$.
Radial integral: Let $u = 1 + r^2 \implies du = 2r \, dr \implies r \, dr = \frac{du}{2}$.
When $r = 0, u = 1$; when $r = 1, u = 2$:
\[
\int_0^1 \frac{r \, dr}{(1+r^2)^2} = \frac{1}{2} \int_1^2 u^{-2} \, du = \frac{1}{2} \left[ -\frac{1}{u} \right]_1^2 = \frac{1}{2}\left( -\frac{1}{2} - (-1) \right) = \frac{1}{2}\left(\frac{1}{2}\right) = \frac{1}{4}
\]
Multiply both factors:
\[
I = \left(\frac{\pi}{2}\right) \left(\frac{1}{4}\right) = \mathbf{\frac{\pi}{8}}
\]

\begin{markingrubric}
\textbf{Mark Distribution:}
$\bullet$ Polar coordinate conversion with Jacobian $r$: \textbf{3 Marks}\\
$\bullet$ Limits of integration for first quadrant: \textbf{2 Marks}\\
$\bullet$ Substitution $u = 1+r^2$ and integration: \textbf{3 Marks}\\
$\bullet$ Final exact answer $\pi/8$: \textbf{2 Marks}
\end{markingrubric}
\end{solutionbox}

% Solution 36
\begin{solutionbox}[Solution to Question 36 (10 Marks)]
\textbf{Step 1: Statement of Stokes' Theorem}
\[
\oint_C \vec{F} \cdot d\vec{r} = \iint_S (\nabla \times \vec{F}) \cdot \hat{n} \, dS
\]
Given $\vec{F} = (2x - y)\hat{i} - y z^2\hat{j} - y^2 z\hat{k}$ over the hemisphere $x^2 + y^2 + z^2 = 1$ ($z \ge 0$).
The boundary curve $C$ is the circle $x^2 + y^2 = 1, z = 0$.

\textbf{Step 2: Line Integral Evaluation ($\oint_C \vec{F} \cdot d\vec{r}$)}
On the boundary circle $C$ in the $xy$-plane: $z = 0, dz = 0$.
Parameterized by $x = \cos\theta, y = \sin\theta$ for $\theta \in [0, 2\pi]$:
\[
dx = -\sin\theta \, d\theta, \quad dy = \cos\theta \, d\theta
\]
Since $z = 0$, the $\hat{j}$ and $\hat{k}$ components of $\vec{F}$ vanish identically!
\[
\vec{F} \cdot d\vec{r} = (2x - y) dx = (2\cos\theta - \sin\theta)(-\sin\theta \, d\theta) = (-2\sin\theta\cos\theta + \sin^2\theta) d\theta
\]
\[
\oint_C \vec{F} \cdot d\vec{r} = \int_0^{2\pi} (-\sin 2\theta + \sin^2\theta) d\theta = 0 + \int_0^{2\pi} \left(\frac{1 - \cos 2\theta}{2}\right) d\theta = \mathbf{\pi}
\]

\textbf{Step 3: Surface Integral Evaluation ($\iint_S (\nabla \times \vec{F}) \cdot \hat{n} \, dS$)}
Compute $\nabla \times \vec{F}$:
\[
\nabla \times \vec{F} = \begin{vmatrix} \hat{i} & \hat{j} & \hat{k} \\ \frac{\partial}{\partial x} & \frac{\partial}{\partial y} & \frac{\partial}{\partial z} \\ 2x - y & -yz^2 & -y^2 z \end{vmatrix}
\]
\begin{itemize}[leftmargin=6mm]
    \item $\hat{i}$-component: $\frac{\partial}{\partial y}(-y^2 z) - \frac{\partial}{\partial z}(-yz^2) = -2yz - (-2yz) = 0$.
    \item $\hat{j}$-component: $\frac{\partial}{\partial z}(2x - y) - \frac{\partial}{\partial x}(-y^2 z) = 0 - 0 = 0$.
    \item $\hat{k}$-component: $\frac{\partial}{\partial x}(-yz^2) - \frac{\partial}{\partial y}(2x - y) = 0 - (-1) = 1$.
\end{itemize}
Thus:
\[
\nabla \times \vec{F} = 0\hat{i} + 0\hat{j} + 1\hat{k} = \mathbf{\hat{k}}
\]
Now by Stokes' Theorem, we can choose the planar base surface $S_1$ (the circular disk $x^2 + y^2 \le 1, z=0$) because both share the identical boundary curve $C$!
On $S_1$, $\hat{n} = \hat{k}$:
\[
(\nabla \times \vec{F}) \cdot \hat{n} = \hat{k} \cdot \hat{k} = 1
\]
\[
\iint_S (\nabla \times \vec{F}) \cdot \hat{n} \, dS = \iint_{S_1} 1 \, dS = \text{Area of unit circle} = \pi (1^2) = \mathbf{\pi}
\]
Since Line Integral ($\pi$) = Surface Integral ($\pi$), **Stokes' Theorem is verified!**

\begin{markingrubric}
\textbf{Mark Distribution:}
$\bullet$ Line integral evaluation yielding $\pi$: \textbf{4 Marks}\\
$\bullet$ Computing $\nabla \times \vec{F} = \hat{k}$: \textbf{3 Marks}\\
$\bullet$ Surface integral evaluation yielding $\pi$: \textbf{2 Marks}\\
$\bullet$ Equating both and concluding verification: \textbf{1 Mark}
\end{markingrubric}
\end{solutionbox}

\newpage
